\section{Введение.}
Данная работа предлагает доказательство гипотезы Бориса Дубровина для особого вида многообразий с коммутативным ассоциативным умножением в касательном пространстве, называемых полиномальными многообразиями Дубровина-Фробениуса. Гипотеза для особенности типа \(A\) заключается в следующем: 
\begin{center}
	\textbf{
		Всякое неприводимое полупростое в общем полиномиальное многообразие Дубровина-Фробениуса, квазиоднородное с градуировкой \(M_{A_\mu}\) изоморфно многообразию Дубровина-Фробениуса для особенности типа \(A\).
	}
\end{center}
Доказательство основывается на приведении алгебр \(T_pM\) к особому виду и представлению их как фактор алгебр по некоторому многочлену и последующей работой с градуировкой и условием интегрируемости для умножения.